\documentclass[11pt]{article}

\usepackage[utf8]{inputenc}
\usepackage[T1]{fontenc}
\usepackage[margin=1in]{geometry}
\usepackage{parskip}
\usepackage{xcolor}
\usepackage{fancyhdr}
\usepackage[expansion=false]{microtype}
\usepackage{hyperref}

\definecolor{accent}{RGB}{91,33,33}
\definecolor{slate}{RGB}{60,60,60}

\hypersetup{
  colorlinks=true,
  linkcolor=accent,
  urlcolor=accent,
  citecolor=accent,
  pdftitle={The Three Days of Darkness: An Analysis},
  pdfauthor={Analysis Document}
}

% Coloured section headings without titlesec/sectsty (base LaTeX only)
\makeatletter
\renewcommand\section{\@startsection{section}{1}{\z@}%
  {-3.5ex \@plus -1ex \@minus -.2ex}{2.3ex \@plus.2ex}%
  {\normalfont\Large\bfseries\color{accent}}}
\renewcommand\subsection{\@startsection{subsection}{2}{\z@}%
  {-3.25ex\@plus -1ex \@minus -.2ex}{1.5ex \@plus .2ex}%
  {\normalfont\large\bfseries\color{accent}}}
\makeatother

\pagestyle{fancy}
\fancyhf{}
\renewcommand{\headrulewidth}{0.4pt}
\lhead{\small\color{slate}The Three Days of Darkness}
\rhead{\small\color{slate}\thepage}

\newcommand{\quoteblock}[2]{%
  \begin{quote}\itshape #1\end{quote}%
  \begin{flushright}\small\color{slate}--- #2\end{flushright}%
}

\title{\color{accent}\bfseries The Three Days of Darkness\\[4pt]
  \large\mdseries A Critical Analysis of the Prophecy and Its\\
  Attribution to Padre Pio and Other Catholic Mystics}
\author{}
\date{\today}

\begin{document}

\maketitle
\thispagestyle{fancy}

\begin{abstract}
\noindent This document examines the popular Catholic prophecy known as the
``Three Days of Darkness''---a predicted period during which the earth is said
to be plunged into a supernatural darkness as a chastisement. The prophecy is
frequently attributed to Saint Pio of Pietrelcina (Padre Pio), to the
contemporary exorcist Father Chad Ripperger, and to a chain of mystics
reaching back through Marie-Julie Jahenny and Blessed Anna Maria Taigi. We
describe the content of the prophecy, trace its historical provenance, weigh
the evidence for and against its authenticity, situate it within Catholic
teaching on private revelation, and conclude with a section on what a believer
might prudently do to prepare---spiritually and practically---without
succumbing to fear or superstition.
\end{abstract}

\tableofcontents
\vspace{1em}

\section{What the Prophecy Predicts}

The ``Three Days of Darkness'' is a prophetic tradition circulating in certain
Catholic devotional circles. In its most common form it predicts that, as a
divine chastisement upon a sinful world, the earth will be enveloped in an
intense, preternatural darkness lasting three days and three nights. The
darkness is described not as a mere astronomical event but as a supernatural
phenomenon accompanied by terror, storms, and the activity of demons.

The single most-quoted formulation of the prophecy is attributed to Blessed
Anna Maria Taigi (1769--1837)~\cite{akin}:

\quoteblock{There shall come over the whole earth an intense darkness lasting
three days and three nights. Nothing can be seen, and the air will be laden
with pestilence which will claim mainly, but not only, the enemies of
religion. It will be impossible to use any man-made lighting during this
darkness, except blessed candles. He who, out of curiosity, opens his window
to look out, or leaves his home, will fall dead on the spot. During these
three days people should remain in their homes, pray the Rosary, and beg God
for mercy.}{attributed to Bl.\ Anna Maria Taigi}

The recurring elements across the various tellings are:

\begin{itemize}
  \item A darkness so complete that no natural or artificial light penetrates
        it---only \emph{blessed} candles (in some versions, only pure beeswax)
        will burn.
  \item A command to stay indoors, seal the windows and doors, and not look
        outside under pain of death.
  \item The air ``infected'' by demons appearing in hideous forms.
  \item A great mortality, said to fall principally upon the impenitent and
        the enemies of the Church.
  \item Continuous prayer---especially the Rosary---as the means of perseverance.
  \item An aftermath in which the world is purified and a period of renewal or
        peace begins.
\end{itemize}

\section{The Chain of Attributions}

\subsection{Blessed Anna Maria Taigi (1769--1837)}
Taigi is the source to whom the classic long quotation is assigned. She was a
Roman wife and mother, beatified in 1920, who was reported to have seen future
events in a luminous orb or ``mystic sun.'' The crucial historical difficulty
is that the specific ``three days'' time-frame does not appear in documents
from her lifetime. It surfaces only around 1863--1864---roughly
twenty-six years after her death---reportedly through her elderly spiritual
director, Raffaele Natali, who recounted it from memory~\cite{ermitano}. The much-cited book
\emph{Private Prophecy} (Italian: \emph{Profezie Private}), said to contain the
text, cannot be located in any major archive.

\subsection{Marie-Julie Jahenny (1850--1941)}
The Breton stigmatist Marie-Julie Jahenny greatly expanded the prophecy in
the late nineteenth and early twentieth centuries. Her reported revelations
add the now-famous detail that only candles of \emph{pure beeswax} will give
light during the darkness~\cite{jahenny}, and they elaborate the instructions to remain
sealed indoors. Her writings are the principal source for the highly specific
practical preparations that circulate today. Her revelations have never
received formal ecclesiastical approval.

\subsection{Padre Pio (St.\ Pio of Pietrelcina, 1887--1968)}
The attribution to Padre Pio is the most popularly repeated and the most
historically problematic. A set of three dramatic ``letters'' describing
apocalyptic visions---darkness, storms, earthquakes, demonic activity---
circulates under his name. However:

\begin{itemize}
  \item The texts have no verifiable manuscript source traceable to Padre Pio.
  \item The Capuchin order, his own religious family and the custodians of his
        writings, has stated that he made no such prediction~\cite{capuchin}.
  \item Padre Pio's authenticated correspondence and spiritual counsel do not
        contain this prophecy.
\end{itemize}

The strong likelihood, therefore, is that the prophecy was attached to his
revered name after the fact to lend it authority.

\subsection{Father Chad Ripperger (b.\ 1964)}
Father Ripperger is a contemporary American priest, exorcist, and author
known for traditional theological and spiritual writing. He has discussed the
Three Days of Darkness in talks and writings as a prophecy held within
traditional devotional circles, treating the question of preparation and the
spiritual dispositions a Catholic should cultivate. His role is best
understood as that of a present-day commentator engaging an existing tradition
rather than an originating visionary; he does not claim a personal private
revelation of the event.

\subsection{Other Names Invoked}
The prophecy is sometimes traced to St.\ Hildegard of Bingen (12th c.),
St.\ Patrick (5th c.), and St.\ Teresa of \'Avila (16th c.). No corresponding
prediction has been found in the authenticated writings of any of these
saints. Their names appear to be invoked to give the prophecy a false
antiquity.

\section{Where It Stands in Catholic Teaching}

Several points are essential for situating this prophecy correctly:

\begin{enumerate}
  \item \textbf{It is not Church teaching.} The phrase ``three days of
        darkness'' does not appear in the Catechism, in conciliar or papal
        teaching, or on the Holy See's official materials. No Catholic is
        bound to believe it.

  \item \textbf{It is private, not public, revelation.} Public revelation
        closed with the death of the last Apostle. Even \emph{approved}
        private revelations (e.g.\ Lourdes, F\'atima) add nothing to the
        deposit of faith and do not oblige belief. The Three Days of Darkness
        has \emph{not} been approved.

  \item \textbf{Canonization does not validate a saint's private visions.}
        That Taigi was beatified or Padre Pio canonized certifies the holiness
        of their lives and (for canonization) attests to miracles---it does
        \emph{not} certify the authenticity of every revelation attributed to
        them.

  \item \textbf{It is absent from Scripture as a future prophecy.} The three
        days of darkness in Exodus 10 were a past plague upon Egypt; the
        darkenings in Revelation are commonly read as referring to other
        realities and do not specify ``three days.''
\end{enumerate}

\section{Assessment of the Evidence}

A sober weighing yields a cautious, skeptical conclusion regarding the
\emph{specific, sensational} form of the prophecy:

\begin{itemize}
  \item The textual chain is weak. The flagship Taigi quotation surfaces
        decades after her death, through a single intermediary, in a book that
        cannot be found.
  \item The Padre Pio attribution is repudiated by his own order.
  \item The appeals to ancient saints are unsupported by their writings.
  \item None of the modern visionaries promoting it have had their revelations
        ecclesiastically approved; at least one prominent promoter has been
        repudiated by his bishop.
\end{itemize}

At the same time, intellectual honesty requires acknowledging what the
tradition gets \emph{right} at the level of principle, independent of the
disputed details: that history can include divine chastisement; that the
Church has always counselled vigilance, repentance, prayer, and reliance on
the sacraments; and that Christians are told to be ready, for they ``know
neither the day nor the hour.'' These are sound whether or not any particular
three-day event ever occurs.

\section{What a Person Can Do to Prepare}

The most defensible posture is to separate \emph{spiritual} preparation, which
is always good regardless of this prophecy's truth, from \emph{material}
preparation, which should be undertaken prudently and without superstition.

\subsection{Spiritual Preparation (always worthwhile)}
\begin{itemize}
  \item \textbf{Live in a state of grace.} Frequent the sacraments---regular
        confession and the Eucharist. This is the Church's perennial counsel
        for being ``ready'' and does not depend on any prophecy.
  \item \textbf{Cultivate a life of prayer,} especially the Rosary, which
        every version of the prophecy commends and which the Church recommends
        universally.
  \item \textbf{Practice trust over fear.} Catholic spirituality consistently
        warns against a fearful, fortune-telling fixation on chastisements.
        The proper response to any prophecy of judgment is conversion and
        confidence in God's mercy, not panic.
  \item \textbf{Examine and amend one's life.} Repentance, charity toward
        others, and detachment from sin are the substance of true readiness.
\end{itemize}

\subsection{Material Preparation (prudent, not superstitious)}
The traditional practical counsel associated with the prophecy~\cite{jahenny,catholicsupply} includes:
\begin{itemize}
  \item Keeping \emph{blessed} candles---pure beeswax in the Jahenny
        tradition---on hand.
  \item Keeping Holy Water and sacramentals (e.g.\ a crucifix, blessed images)
        in the home.
  \item Staying indoors, with windows and doors closed, and praying through
        the event.
\end{itemize}

These can be adopted as ordinary acts of devotion and reasonable household
readiness. They overlap, in fact, with sensible general emergency
preparedness---candles, water, food, and a calm plan to shelter in place---
which is prudent for any number of mundane emergencies (storms, outages)
quite apart from any prophecy. The key caution is theological: a blessed
candle or Holy Water is a \emph{sacramental}, not a magical talisman. Its
power, in Catholic understanding, lies in the prayer of the Church and the
disposition of faith, not in the object as such. Treating these items as
mechanical guarantees of survival would be superstition, which the Church
condemns~\cite{ccc}.

\subsection{A Balanced Bottom Line}
A reasonable Catholic response is therefore: pursue the spiritual preparation
wholeheartedly (it is simply the Christian life), keep ordinary sacramentals
and sensible emergency supplies without anxiety, withhold firm belief from the
unverified specifics, and refuse to let any unapproved prophecy become a
source of fear, financial exploitation, or distraction from the Gospel.

\section{Conclusion}

The Three Days of Darkness is a vivid and widely circulated prophecy whose
\emph{specific} form rests on a fragile evidentiary chain: a flagship
quotation that appears decades posthumously through a single witness, an
attribution to Padre Pio that his own order denies, and appeals to ancient
saints unsupported by their writings. It is not Church teaching, it is not in
Scripture as a future event, and no approved apparition confirms it. Catholics
are entirely free to disbelieve the particulars.

Yet the \emph{dispositions} it calls for---repentance, prayer, the
sacraments, trust in God, and quiet readiness---are simply the Christian life,
and these lose nothing if the prophecy is false. The wise course is to embrace
that interior preparation, adopt only prudent and non-superstitious external
measures, and meet all such predictions with the Gospel's own counsel: watch,
pray, do not be afraid.

\begin{thebibliography}{9}

\bibitem{akin}
Jimmy Akin,
\emph{Will There Be Three Days of Darkness?},
Catholic Answers Magazine (online edition), 18 September 2025.
\url{https://www.catholic.com/magazine/online-edition/will-there-be-three-days-of-darkness}

\bibitem{ermitano}
\emph{El Ermitaño}, No.\ 155 (26 October 1871), reporting the c.\ August 1864
recollections of Raffaele Natali, spiritual director of Bl.\ Anna Maria Taigi.
(Cited in Akin, ref.\ \cite{akin}.)

\bibitem{ccc}
Catholic Church,
\emph{Catechism of the Catholic Church}, 2nd ed.,
on private revelation (\S 65--67) and on superstition (\S 2110--2117).
Libreria Editrice Vaticana.

\bibitem{jahenny}
Revelations attributed to Marie-Julie Jahenny (1850--1941), the ``Breton
Stigmatist,'' as summarized in popular devotional literature concerning the
Three Days of Darkness and the use of blessed pure-beeswax candles.
(Unapproved private revelation; see discussion and caveats in Akin, ref.\ \cite{akin}.)

\bibitem{capuchin}
Statement of the Capuchin Order denying that St.\ Pio of Pietrelcina (Padre
Pio) authored the apocalyptic ``letters'' predicting the Three Days of
Darkness. (Reported in Akin, ref.\ \cite{akin}.)

\bibitem{catholicsupply}
Catholic Supply,
\emph{Preparing for the 3 Days of Darkness} (product and preparation guide).
\url{https://shop.catholicsupply.com/preparing-for-the-3-days-of-darkness-catholic-supply.aspx}
(Cited as an example of the practical/material counsel circulating in devotional commerce, not as an endorsement.)

\bibitem{scripture}
\emph{The Holy Bible}: Exodus 10:21--29 (the ninth plague upon Egypt) and
Revelation 6:12, 16:10; with the Gospel counsel to watchfulness,
Matthew 24:42 and 25:13 (``you know neither the day nor the hour'').

\end{thebibliography}

\end{document}
